\structuralelement{Заключение}

В ходе выполнения выпускной квалификационной работы поставленная цель
достигнута: разработано программное обеспечение автономного мобильного
омни-робота, обеспечивающее сопровождение пользователя в квартирной
среде с динамическими препятствиями на базе стохастической реализации
модельного прогнозирующего управления.

В рамках выполнения задач ВКР получены следующие основные результаты.
\begin{enumerate}
    \item Проведён анализ существующих подходов к локальному планированию
          траектории и методов SLAM, подтверждена целесообразность выбора
          MPPI как основного локального контроллера и slam\_toolbox как
          решения задачи SLAM на базе 2D-лидара; в качестве
          интегрирующего программного стека обоснованно выбран Nav2.
    \item Построена математическая модель омни-платформы с колёсами
          Mecanum: в обозначениях кинематической схемы введены связные
          $L$, $l$, $\psi$, $r$, $\dot{\varphi}_i$, выведены
          кинематические соотношения в непрерывной и дискретной формах,
          записаны формулы связи обобщённых скоростей корпуса и угловых
          скоростей колёс, зафиксированы ограничения управления и
          условия безопасности относительно препятствий.
    \item Формализован принцип модельного прогнозирующего управления и
          его стохастической реализации MPPI; сформирована составная
          функция потерь, отражающая одновременные требования по
          удержанию дистанции до цели, центрированию цели в поле зрения
          камеры, безопасному обходу препятствий, плавности управления и
          остановке при потере цели; показана эквивалентность
          построенной функции потерь совокупности критиков MPPI.
    \item Разработана воспроизводимая контейнерная среда исполнения на
          базе Docker, ROS~2 Humble и Gazebo, с автоматизированной
          оркестрацией стартовой последовательности через tmux и
          удалённым доступом к графическому интерфейсу через noVNC.
    \item Реализована модель омни-робота с 2D-лидаром и фронтальной
          RGB-камерой в формате URDF/xacro; развёрнут набор прикладных
          ROS-узлов (мосты тем и преобразований, детектор цели, мост
          действий навигации, контроллер подвижных сущностей).
    \item Реализован пользовательский критик \texttt{TargetCameraCritic}
          в составе \texttt{nav2\_mppi\_controller}, непосредственно
          реализующий слагаемое функции потерь, отвечающее за
          центрирование цели в поле зрения камеры; критик корректно
          отключается при потере наблюдаемости цели, что соответствует
          теоретической формулировке.
    \item Реализован пользовательский критик
          \texttt{PotentialFieldCritic}, формирующий направленный
          (анизотропный) штраф вдоль градиента поля препятствий
          (слагаемое $\Psi_{\text{pf},k}$ функции потерь) с
          активацией по порогу costmap и автоматическим выключением
          вблизи финальной позы; критик восполняет ограничение
          изотропных штрафов \texttt{ObstaclesCritic}/\texttt{CostCritic}
          для голономной омни-платформы.
\end{enumerate}

Таким образом, поставленная задача разработки программно-алгоритмического
комплекса для омни-робота в задаче сопровождения пользователя в
квартирной среде решена в объёме, достаточном для демонстрации всех
требуемых функциональных свойств системы. Получена устойчивая и
модульная архитектура, пригодная для дальнейшего расширения: замены
цветового детектора на более общий алгоритм распознавания человека,
перехода от симуляции к реальной платформе, настройки весов функции
потерь под конкретные сценарии эксплуатации.

Перспективы дальнейшей разработки темы включают:
\begin{itemize}
    \item замену модельного детектора цели на детектор на основе
          современных моделей машинного обучения (например, YOLO-класса)
          для работы с произвольным пользователем в бытовой среде;
    \item проведение прямого A/B-сравнения по тем же метрикам
          с отключённым \texttt{PotentialFieldCritic} для строгой
          количественной оценки его вклада в безопасность;
    \item переход от симуляционной валидации к экспериментам на
          реальной платформе Mecanum-типа с последующим
          переопределением весов функции потерь по данным полевых
          испытаний.
\end{itemize}
